\section{Cut-based recourse: exact output, concave--convex residuals and proofs}\label{app:recourse-cuts}

This appendix proves Theorems~\ref{thm:cr-oracle} and~\ref{thm:cr-search}
and Proposition~\ref{prop:cr-growth} of Section~\ref{sec:cuts}, the
exact-output theorem for a cut residual (Theorem~\ref{thm:cr-exact}), the
claims of Remark~\ref{rem:cr-mixed} on concave--convex residuals, and the
results deferred from Section~\ref{sec:balanced}. The first three
subsections use the notation of Section~\ref{sec:cuts}:
$F(x)=\frac12x^{\T}Hx+b^{\T}x+c$ is a rational quadratic, $\mathcal K$ is a
core of $k$ continuous coordinates with domain $[0,1]^k$, $\mathcal R$ is the
residual, and $\alpha_i$, $\delta_i$ and $y(z)$ are as defined before
Proposition~\ref{prop:cr-cut}; the last uses that of
Section~\ref{sec:balanced}.

\subsection{Proofs for Section~\ref{sec:cuts}}\label{app:cr-proofs}

The proofs below use the bag-cell error $e_B(C)$ of Section~\ref{sec:local}
with $B=\mathcal K$, and $e_j$, $\beta_V(C)$, $\mathcal Q_j$,
$\mathcal Q_j'$, $\lambda_j$ and $U$ as in CORE (Algorithm~\ref{alg:core}).

\begin{proof}[Proof of Theorem~\ref{thm:cr-oracle}]
By Lemma~\ref{lem:cr-endpoint} the conditional minimum is attained at some
$y(z)$, and by Proposition~\ref{prop:cr-cut} minimizing over $z$ is a
minimum cut problem. Every feasible flow has value at most the capacity of
every cut (weak duality), so a flow of value $\operatorname{cap}(S_z)$
proves that $S_z$ is a minimum cut; checking it needs the capacity bounds,
conservation at each residual node and one comparison. A maximum flow and a
minimum cut of equal value exist and are found, for instance, by the
Edmonds--Karp algorithm \cite{EdmondsKarp1972}, whose number of arithmetic
operations is polynomial in the number of nodes and arcs, here
$\abs{\mathcal R}+2$ and at most $2\abs{\mathcal R}$ plus twice the number
of nonzero residual couplings. After multiplication by a common denominator
of polynomial bit length, the capacities are integers, every augmentation is
by an integer, and every flow value is an integer at most the total
capacity, so all intermediate numbers have polynomial bit length. The
coefficients of Proposition~\ref{prop:cr-cut}(b) are obtained by exact
rational evaluation.
\end{proof}

\begin{proof}[Proof of Theorem~\ref{thm:cr-search}]
Let $C$ be a cell of side $h_j$ and $x\in X$ with $x_{\mathcal K}\in C$.
Lemma~\ref{lem:cr-cell} with $B=\mathcal K$ gives
$F(x)\ge\min\{V(v):v\text{ a corner of }C\}-e_{\mathcal K}(C)\ge\beta_V(C)$,
because the core coordinates are continuous and $e_{\mathcal K}(C)\le e_j$.

(a) The unit cube contains $x^*_{\mathcal K}$. If $C\in\mathcal Q_j$ contains
it, then $\beta_V(C)\le F(x^*)=\OPT\le U$, so $C$ is retained and every
child of $C$ that contains $x^*_{\mathcal K}$ belongs to $\mathcal Q_{j+1}$.
(b), (c) These are Theorem~\ref{thm:cr-filter} with $B=\mathcal K$, the
cells $\mathcal Q_j$, $e=e_j$, $\underline V=V$ and $\varepsilon_{\mathrm{or}}=0$; its hypothesis
that some cell of $\mathcal Q_j$ contains $x^*_{\mathcal K}$ holds by (a), and
$\lambda_j\le\OPT$ by the first paragraph applied to $x^*$.
(d) Children of retained cells cover them, so the cells of $\mathcal Q_j$
together with the cells discarded at levels $i<j$ cover $[0,1]^k$. If
$x_{\mathcal K}$ lies in a cell of $\mathcal Q_j$, then $F(x)\ge\lambda_j$ by
the first paragraph. If it lies in a cell $C$ discarded at level $i$, then
$F(x)\ge\beta_V(C)>U^{(i)}\ge U$, where $U^{(i)}$ is the incumbent value at
level $i$. Each step uses only recorded values.
\end{proof}

\begin{proof}[Proof of Proposition~\ref{prop:cr-growth}]
Let $N_j$ be the number of points $v\in h_j\Z^k\cap[0,1]^k$ with
$V(v)\le\OPT+2e_j$. Level $0$ makes $2^k$ queries. By
Theorem~\ref{thm:cr-search}(c), every retained level-$j$ cell has a corner
counted in $N_j$, and a point is a corner of at most $2^k$ cells, so
$\abs{\mathcal Q_j'}\le2^kN_j$. Each retained cell has $2^k$ children with
$2^k$ corners each, so level $j+1$ makes at most $8^kN_j$ queries. A point
counted in $N_j$ satisfies $g\norm{v-v^*}^2\le2e_j=kLh_j^2/4$, so each of
its coordinates lies within $\frac{h_j}2\sqrt{k\kappa_{\mathcal K}}$ of that
of $v^*$; an interval of length $h_j\sqrt{k\kappa_{\mathcal K}}$ contains at
most $1+\sqrt{k\kappa_{\mathcal K}}$ points of $h_j\Z$. Hence
$N_j\le(1+\sqrt{k\kappa_{\mathcal K}})^k$; sum over $j<J$.
\end{proof}

\subsection{Exact output with a cut residual}

Let $\Delta_{\mathcal K}=\Lambda_c\Lambda_e^2$, where $\Lambda_c$ is the
least positive common denominator of the monomial coefficients of $F$ and
$\Lambda_e$ that of the residual bounds after inward rounding; both have
$O(I)$ bits. For a residual
endpoint vector $y$ put $F_y(v)=F(v,y)$ and
$\Psi(y)=\min_{v\in[0,1]^k}F_y(v)$.

\begin{lemma}[Uniform height over residual endpoint vectors]\label{lem:cr-height}
Assume (R1). Let $\Omega_{\mathcal K}$ be the constant $\Omega$
of~\eqref{eq:exact-constants} computed for $F_y$ on $[0,1]^k$ with the
common denominator $\Delta=\Delta_{\mathcal K}$. Then $\Omega_{\mathcal K}$
does not depend on $y$, $\log\Omega_{\mathcal K}=O(I^2)$, and for every
residual endpoint vector $y$ the value $\Psi(y)$ is a rational with reduced
denominator at most $\Omega_{\mathcal K}$; so is $\OPT$.
\end{lemma}

\begin{proof}
The monomial coefficients of $F_y$ are $\frac12H_{ii}$ for $i\in\mathcal K$,
$H_{ij}$ for $i<j$ in $\mathcal K$, the linear coefficients
$b_i+\sum_{j\in\mathcal R}H_{ij}y_j$ and the constant
$c+\sum_{j\in\mathcal R}(b_jy_j+\frac12H_{jj}y_j^2)
+\sum_{j<l,\,j,l\in\mathcal R}H_{jl}y_jy_l$. Each is a sum of products of a
coefficient of $F$ with at most two residual bounds, so $\Delta_{\mathcal K}$
is a common denominator of all of them and of the endpoints $0,1$. The
Hessian of $F_y$ is $H_{\mathcal K\mathcal K}$ for every $y$, and the
constant $\Omega$ of~\eqref{eq:exact-constants} depends only on $\Delta$ and
the Hessian, so $\Omega_{\mathcal K}$ does not depend on $y$; its bit
length is $O(kI)=O(I^2)$, as in Section~\ref{sec:exact-data}.
Corollary~\ref{cor:height}(b), applied to the box QP $F_y$ on
$[0,1]^k$, bounds the reduced denominator of $\Psi(y)$ by
$\Omega_{\mathcal K}$. By Lemma~\ref{lem:cr-endpoint},
$\OPT=\min_v\min_{y}F(v,y)$ with $y$ ranging over residual endpoint vectors,
so $\OPT=\Psi(y^*)$ for some such $y^*$.
\end{proof}

The square of $\Lambda_e$ is needed: the term $\frac12y_iy_j$ at residual
endpoints $y_i=y_j=\frac12$ equals $\frac18$, while the data have common
denominator $2$.

\begin{theorem}[Exact output with a cut residual]\label{thm:cr-exact}
Assume~\eqref{eq:cr-class} with continuous core $[0,1]^k$ and rational data,
and run CORE (Algorithm~\ref{alg:core}) with the oracle of
Theorem~\ref{thm:cr-oracle}, $L=\max_{i\in\mathcal K}\max\{H_{ii},0\}$,
accuracy $\varepsilon=\Omega_{\mathcal K}^{-2}/2$ and level limit $J$ the
first level with $e_J\le\varepsilon$. Let $y_U$ be the residual part of the
final incumbent and $\beta=\min\{\lambda_j,U\}$ the returned bound.
Minimize $F_{y_U}$ exactly over $[0,1]^k$ by enumerating the $3^k$ faces of
the cube, solving the stationarity system on each face whose free Hessian
block is nonsingular, and keeping the solutions that lie in the face. The best candidate $\hat v$ makes $(\hat v,y_U)$ a global
minimizer, and $(\hat v,y_U)$ passes the test of
Proposition~\ref{prop:accept} with the denominator bound
$\Omega_{\mathcal K}$ and this $\beta$. A checker needs only
the trace of CORE, $\Omega_{\mathcal K}$, and the feasibility, exact value
and reduced denominator of the output. If in addition $V$ has core growth
with constant $g$ as in Proposition~\ref{prop:cr-growth}, the total bit
complexity is $8^k(1+\sqrt{k\kappa_{\mathcal K}})^k\poly(I)$ with an
absolute exponent, and $g$ is not an input.
\end{theorem}

\begin{proof}
By Theorem~\ref{thm:cr-search}(b), CORE stops at level $J$ at the latest,
with $U-\beta\le\varepsilon<\Omega_{\mathcal K}^{-2}$. Apply
Lemma~\ref{lem:statpoly} to the box QP $F_{y_U}$ on $[0,1]^k$: let $v^\circ$
be a minimizer and $v$ a vertex of its stationary polytope. By part~(b), $v$ is a
minimizer, and by part~(c) the Hessian block of $F_{y_U}$ on the set
$J_0(v)$ of coordinates of $v$ in $(0,1)$ is positive definite. So $v$ is the unique solution of the stationarity
system on the face of the cube that fixes the coordinates outside $J_0(v)$
at their values in $v$, and the enumeration contains it. All candidates lie
in $[0,1]^k$, so the best one attains $\Psi(y_U)$.

The oracle returns endpoint vectors, so $y_U$ is one. With $(v_U,y_U)$ the
final incumbent, $F(\hat v,y_U)=\Psi(y_U)\le F(v_U,y_U)=U$, and
$\beta\le\OPT$ by Theorem~\ref{thm:cr-search}(d). By
Lemma~\ref{lem:cr-height}, $\OPT$ has denominator at most
$\Omega_{\mathcal K}$ and $F(\hat v,y_U)=a/W$ in lowest terms with
$W\le\Omega_{\mathcal K}$. Hence
$F(\hat v,y_U)-\beta\le U-\beta<\Omega_{\mathcal K}^{-2}\le
1/(\Omega_{\mathcal K}W)$, and Proposition~\ref{prop:accept} shows that
$(\hat v,y_U)$ is a global minimizer. The checker verifies
$\beta\le\OPT$ by Theorem~\ref{thm:cr-search}(d) and then the inequality
$F(\hat v,y_U)-\beta<1/(\Omega_{\mathcal K}W)$, which proves optimality by
Proposition~\ref{prop:accept}.

For the complexity, if $L=0$ then $J=0$. Otherwise
$\log(kL)=O(I)$ and $\log\Omega_{\mathcal K}=O(I^2)$, so
$J=O(\log\Omega_{\mathcal K}+\log(kL)+1)=O(I^2)$.
Under core growth, Proposition~\ref{prop:cr-growth} bounds the queries by
$2^k+8^kJ(1+\sqrt{k\kappa_{\mathcal K}})^k$. Each query is a cut problem
whose data have bit length polynomial in $I+kJ$, hence $\poly(I)$ by
Theorem~\ref{thm:cr-oracle}, with an absolute exponent. The face
enumeration costs $3^k\poly(I)$, and $2^k+3^k\le2\cdot8^k$.
\end{proof}

Without core growth the procedure still terminates, after at most
$J=O(I^2)$ levels, but Example~\ref{ex:cr-flat} shows that the number of
queries can then be $2^{\Omega(kJ)}$.

\subsection{Concave--convex residuals}

This subsection proves the claims of Remark~\ref{rem:cr-mixed}. Fix a
rational core point $v$. Let $\mathcal R=\mathcal R_-\cup\mathcal R_+$ with
$H_{ii}\le0$ for $i\in\mathcal R_-$, signs $o$ with $o_io_jH_{ij}\le0$ for
all $i\ne j$ in $\mathcal R$, continuous coordinates in $\mathcal R_+$ and
$H_{\mathcal R_+\mathcal R_+}\succeq0$. For $S\subseteq\mathcal R_-$
let $y_-(S)\in\R^{\mathcal R_-}$ have coordinates
$\alpha_i+\delta_i[i\in S]$, and put
\[
\Phi(S)=\min\Bigl\{F\bigl(v,y_-(S),y_+\bigr):
y_+\in\prod_{i\in\mathcal R_+}[\ell_i,u_i]\Bigr\}.
\]
With $\mathcal R_+=\emptyset$ this is the class~\eqref{eq:cr-class}; with
$\mathcal R_-=\emptyset$ the residual problem is convex.

\begin{proposition}[Submodular conditional value]\label{prop:cr-submod}
(a) The minimum of $F(v,\cdot)$ over the continuous residual box and over
the mixed residual domain equals $\min_{S\subseteq\mathcal R_-}\Phi(S)$.
(b) Each $\Phi(S)$ is the value of a rational convex box QP.
(c) For all $S,S'\subseteq\mathcal R_-$,
$\Phi(S)+\Phi(S')\ge\Phi(S\cap S')+\Phi(S\cup S')$.
\end{proposition}

\begin{proof}
(a) As in Lemma~\ref{lem:cr-endpoint}, move each coordinate of
$\mathcal R_-$ to an endpoint without increasing $F$; the resulting point is
in the mixed domain because $\mathcal R_+$ is continuous. (b) After the
coordinates in $\mathcal R_-$ are fixed, the objective in $y_+$ has Hessian
$H_{\mathcal R_+\mathcal R_+}\succeq0$.
(c) Substitute out the coordinates with $\delta_i=0$ and put
$y_i=\alpha_i+\delta_it_i$ with $t\in[0,1]^{\mathcal R}$. The transformed
objective $\tilde F$ has off-diagonal Hessian entries
$\delta_i\delta_jH_{ij}=o_io_j(u_i-\ell_i)(u_j-\ell_j)H_{ij}\le0$. For
$t,t'\in[0,1]^{\mathcal R}$,
$\tilde F(t)+\tilde F(t')\ge\tilde F(t\wedge t')+\tilde F(t\vee t')$: the
terms in one coordinate cancel, because $\{t_i,t'_i\}=\{(t\wedge
t')_i,(t\vee t')_i\}$; for a pair $i,j$ the two sides agree if $t_i-t'_i$ and
$t_j-t'_j$ do not have opposite signs, and otherwise
$t_it_j+t'_it'_j-(t\wedge t')_i(t\wedge t')_j-(t\vee t')_i(t\vee t')_j
=(t_i-t'_i)(t_j-t'_j)\le0$, which, multiplied by a nonpositive coefficient,
is nonnegative. Let $t^S$ and $t^{S'}$ attain $\Phi(S)$ and $\Phi(S')$ in
the transformed coordinates of $\mathcal R_+$. Then
$(\mathbf 1_S,t^S)\wedge(\mathbf 1_{S'},t^{S'})
=(\mathbf 1_{S\cap S'},t^S\wedge t^{S'})$, similarly for $\vee$ with
$S\cup S'$, and both completions are feasible. Hence
$\Phi(S)+\Phi(S')\ge\tilde F(\mathbf 1_{S\cap S'},t^S\wedge t^{S'})
+\tilde F(\mathbf 1_{S\cup S'},t^S\vee t^{S'})\ge\Phi(S\cap S')+\Phi(S\cup S')$.
\end{proof}

Positive semidefiniteness is not used in (c); it makes each $\Phi(S)$
computable exactly in polynomial time \cite{KozlovTarasovKhachiyan1980}.
Let $m=\abs{\mathcal R_-}$ and $\hat\Phi(S)=\Phi(S)-\Phi(\emptyset)$. For an
ordering $\pi$ of $\mathcal R_-$ with prefixes
$S^\pi_0=\emptyset,S^\pi_1,\dots,S^\pi_m=\mathcal R_-$, the \emph{greedy
vector} $a^\pi\in\R^{\mathcal R_-}$ has
$a^\pi_{\pi(l)}=\hat\Phi(S^\pi_l)-\hat\Phi(S^\pi_{l-1})$. We write
$a(S)=\sum_{i\in S}a_i$.

\begin{proposition}[Greedy-base certificate]\label{prop:cr-greedy}
(a) For every ordering $\pi$ and every $S\subseteq\mathcal R_-$,
$a^\pi(S)\le\hat\Phi(S)$, and $a^\pi(\mathcal R_-)=\hat\Phi(\mathcal R_-)$.
(b) If $\bar a=\sum_\pi\lambda_\pi a^\pi$ is a convex combination, then
$\min_yF(v,y)\ge\Phi(\emptyset)+\sum_{i\in\mathcal R_-}\min\{0,\bar a_i\}$,
the minimum taken over the residual domain.
(c) Some convex combination of at most $m+1$ greedy vectors attains
equality in (b) \cite{Edmonds1970}. Such a combination, a minimizing set
$S^*$ and exact minimizers of the convex box QPs involved can be computed in
time polynomial in the input length.
(d) Given the orderings, the weights, the $O(m^2)$ values $\Phi(S^\pi_l)$
with minimizers, and a set $S^*$ with its completion, a checker verifies the
bound in (b) and its attainment by $S^*$ with polynomially many rational
operations and no optimization. It checks a value $\Phi(S)$ by feasibility,
exact evaluation and the sign conditions $\partial_iF=0$ at interior
coordinates, $\partial_iF\ge0$ at lower bounds and $\partial_iF\le0$ at upper
bounds for $i\in\mathcal R_+$, which suffice by convexity once
$H_{\mathcal R_+\mathcal R_+}\succeq0$ and the sign conditions of the class
have been verified exactly.
\end{proposition}

\begin{proof}
(a) Let $i=\pi(l)\in S$, and let $S_{<i}$ be the set of elements of $S$
that precede $i$ in $\pi$. Since $S_{<i}\subseteq S^\pi_{l-1}$ and
$i\notin S^\pi_{l-1}$, submodularity (Proposition~\ref{prop:cr-submod}(c))
applied to $S_{<i}\cup\{i\}$ and $S^\pi_{l-1}$ gives
$a^\pi_i=\hat\Phi(S^\pi_l)-\hat\Phi(S^\pi_{l-1})
\le\hat\Phi(S_{<i}\cup\{i\})-\hat\Phi(S_{<i})$. Summing over $i\in S$
telescopes to $a^\pi(S)\le\hat\Phi(S)$, because $\hat\Phi(\emptyset)=0$;
the sum over all positions is $\hat\Phi(\mathcal R_-)$.

(b) By (a) and convexity, $\bar a(S)\le\hat\Phi(S)$ for every $S$, so
$\Phi(S)=\Phi(\emptyset)+\hat\Phi(S)\ge\Phi(\emptyset)+\bar a(S)
\ge\Phi(\emptyset)+\sum_i\min\{0,\bar a_i\}$; conclude with
Proposition~\ref{prop:cr-submod}(a).

(c) Edmonds' theorem gives
$\min_S\hat\Phi(S)=\max\{\sum_i\min\{0,\bar a_i\}:\bar a\in B(\hat\Phi)\}$,
where the base polytope $B(\hat\Phi)$ has the greedy vectors as its
vertices; by Carath\'eodory's theorem, the maximum of this concave
piecewise-linear function is attained at a combination of at most $m+1$ of
them. For
computation, consider the linear program
$\max\{t:\ 0\le\xi\le\mathbf 1,\ t+(a^\pi)^{\T}\xi\le0\ \text{for all }\pi\}$,
whose dual is $\min\{\sum_ir_i:\ \sum_\pi\lambda_\pi a^\pi+r\ge0,\
\sum_\pi\lambda_\pi=1,\ \lambda,r\ge0\}$. For $\xi\in[0,1]^m$, an ordering
$\pi$ with $\xi_{\pi(1)}\ge\dots\ge\xi_{\pi(m)}$ maximizes $(a^\pi)^{\T}\xi$
over all orderings \cite{Edmonds1970}, and with $\xi_{\pi(m+1)}:=0$
\[
(a^\pi)^{\T}\xi=\sum_{l=1}^m\bigl(\xi_{\pi(l)}-\xi_{\pi(l+1)}\bigr)
\hat\Phi(S^\pi_l)+\bigl(1-\xi_{\pi(1)}\bigr)\hat\Phi(\emptyset)
\]
is a convex combination of values of $\hat\Phi$, hence at least
$\min\hat\Phi$; equality holds for $\xi=\mathbf 1_S$ with $S$ minimizing
$\hat\Phi$. So the optimal value of the program is $-\min\hat\Phi$, a
separation call costs $m+1$ convex-QP evaluations, and for an optimal $\xi$
every prefix with positive weight in this combination is a minimizing set.
Each $\Phi(S)$ is the optimal value of a convex box QP with Hessian
$H_{\mathcal R_+\mathcal R_+}$ and has polynomial encoding length by
Corollary~\ref{cor:height}(b); hence the polyhedron has polynomial facet
complexity, and the ellipsoid method solves the program and finds an optimal
basic dual solution in polynomial time
\cite[Theorem~6.4.9, Lemma~6.5.15]{GrotschelLovaszSchrijver1988}. A basic
dual solution has at most $m+1$ positive $\lambda_\pi$, and at optimality
$r_i=\max\{0,-\bar a_i\}$ with $\bar a=\sum_\pi\lambda_\pi a^\pi$, so
$\sum_i\min\{0,\bar a_i\}=\min\hat\Phi$.

(d) follows from (a), (b) and the optimality conditions for convex problems
on a box.
\end{proof}

\subsection{Proofs for Section~\ref{sec:balanced}}

\begin{proof}[Proof of Lemma~\ref{lem:chaincut}]
List $G_i$ as $a_{i,0},\dots,a_{i,k_i}$, in increasing order if $o_i=1$ and
in decreasing order if $o_i=-1$, so that $\delta_{i,l}=a_{i,l}-a_{i,l-1}$
has the sign of $o_i$. Encode $y_i=a_{i,l}$ by $z_i\in\{0,1\}^{k_i}$ with
$z_{i,l'}=1$ exactly for $l'\le l$. This is a bijection between $G_i$ and
the nonincreasing vectors $z_{i,1}\ge\dots\ge z_{i,k_i}$, and on these
vectors
\[
y_i=a_{i,0}+\sum_l\delta_{i,l}z_{i,l},\qquad
\varphi(y_i)=\varphi(a_{i,0})
+\sum_l\bigl[\varphi(a_{i,l})-\varphi(a_{i,l-1})\bigr]z_{i,l}
\]
for every function $\varphi$ on $G_i$. Apply the second formula to
$\varphi_i(t)=\frac12H_{ii}t^2+b_it+\chi_i(t)$, and substitute the first into
each cross term $H_{ij}y_iy_j$, $i<j$. On nonincreasing vectors this writes
$F_\chi$ as a function $\Phi(z)$ of the binary vector $z=(z_{i,l})$ of the
form of Proposition~\ref{prop:cr-cut}(a), after the pairs $(i,l)$ are
numbered, with pair coefficients
$H_{ij}\delta_{i,l}\delta_{j,l'}\le0$ for $i\ne j$ and none within one
coordinate. Take the network of Proposition~\ref{prop:cr-cut}(a) for
$\Phi$ and add an arc $(i,l+1)\to(i,l)$ of infinite capacity for each $i$ and
$1\le l<k_i$. A cut $S_z$ has finite capacity exactly when no such arc
leaves it, that is, when $z$ is nonincreasing in each coordinate; its
capacity is then given by~\eqref{eq:cr-cutid}. The all-zero vector is
nonincreasing, so a minimum cut is finite and decodes to a minimizer of
$F_\chi$, and a flow of equal value certifies it by weak duality. Fixing
$y_i=v$ is the same construction with $G_i=\{v\}$, and the restricted
function is still balanced.
\end{proof}

\begin{proof}[Proof of the work bounds in Theorem~\ref{thm:balanced}(b)]
If $L=0$, then $P=\emptyset$ and CT (Algorithm~\ref{alg:ct}) makes one cut
computation. Otherwise point growth with constant $g$ gives weighted growth
with $\gamma=g/L$, because $\norm d_{\mathbf L}^2\le L\norm d^2$, so
$\bar\kappa\le\kappa$. By the proof of Theorem~\ref{thm:approx}(b), CT
succeeds at the latest at the trial $\mu^*$ with
$2^{\mu^*}\le6\sqrt{\bar\kappa}\le6\sqrt\kappa$, so at most
$\log_2(6\sqrt\kappa)$ trials are run; the trial with $\theta=2^{-\mu}$ has
at most $K_\mu\le K_{\mu^*}=O(\sqrt\kappa\log n)$ nodes per coordinate, and
each trial has $O(I+q+1)$ stages. A stage performs at most $1+nK_\mu$ cut
computations on networks with at most $nK_\mu+2$ nodes and
$3nK_\mu+n^2K_\mu^2$ arcs. As in
Appendix~\ref{app:growth}, grid nodes have denominators dividing
$\Gamma_X2^\alpha$ and the values of $F$ and of the corrections at grid
points have denominators dividing $8\Gamma_F\Gamma_X^24^\alpha$, with
$\alpha=O(I+q+\mu K_\mu)$. Each capacity of Lemma~\ref{lem:chaincut} with
$\chi_i=-d_i$ is a sum of $O(nK_\mu)$ terms, each a coefficient of $F$ times
at most two grid nodes or differences of grid nodes, a correction value, or
half of such a term. So the capacities have a common denominator dividing
$8\Gamma_F\Gamma_X^24^\alpha$ and numerators of $O(I+\alpha+\log(nK_\mu))$
bits. Give each arc of infinite capacity the capacity one plus the sum of
all finite capacities; this changes neither the minimum cuts nor the
maximum-flow value. After scaling by the common denominator, the
Edmonds--Karp algorithm performs a number of augmentations polynomial in the
network size, each on integers bounded by the total capacity. Multiplying
these polynomial bounds gives the bound for CT. Under point growth, the
proof of Theorem~\ref{thm:exact} shows that EX (Algorithm~\ref{alg:ex})
stops at an accuracy $q=\poly(I)+O(\log\kappa)$ after $O(\log q)$ calls of
CT, and that REC then needs only Gaussian elimination; this gives the bound
for EX.
\end{proof}
